\begin{figure}[ht]
		\centering
		\begin{tikzpicture}[>=Stealth, line width=0.8pt, font=\small,
			nodo/.style={circle, fill=white, draw=black,
				inner sep=0pt, minimum size=5pt, line width=0.8pt},
			tacca/.style={thin, gray!70}]
			
			\def\a{2.0}     % unità di lunghezza (= a)
			\def\tw{0.22}
			\def\thA{0.38}
			\def\thR{0.32}
			\def\bw{0.16}
			\def\gw{0.30}
			\def\sr{0.08}
			
			%% trave
			\draw[line width=2.2pt] (0,0) -- ({2*\a},0);
			
			%% etichette nodi
			\node[above=2pt] at (0,0)      {$A$};
			\node[above=2pt] at (\a,0)     {$B$};
			\node[above=2pt] at ({2*\a},0) {$C$};
			
			%% cerniera in A
			\draw (0,0) -- ++(-\tw,-\thA) -- ++(2*\tw,0) -- cycle;
			\draw (-\gw,-\thA) -- (\gw,-\thA);
			\foreach \x in {-0.24,-0.12,0,0.12,0.24}
			\draw[thin] (\x,-\thA) -- ++(-0.09,-0.11);
			\node[nodo] at (0,0) {};
			
			%% carrello intermedio in B
			\draw (\a,0) -- ++(-\tw,-\thR) -- ++(2*\tw,0) -- cycle;
			\draw (\a-\gw,-\thR) -- (\a+\gw,-\thR);
			\draw (\a-\bw,-\thR-0.12) circle (2.4pt);
			\draw (\a+\bw,-\thR-0.12) circle (2.4pt);
			\draw (\a-\gw,-\thR-0.24) -- (\a+\gw,-\thR-0.24);
			\foreach \x in {-0.24,-0.12,0,0.12,0.24}
			\draw[thin] (\a+\x,-\thR-0.24) -- ++(-0.09,-0.11);
			\fill[white] (\a-\sr,0) arc(180:360:\sr) -- cycle;
			\draw (\a-\sr,0) arc(180:360:\sr);
			
			%% carrello terminale in C
			\draw ({2*\a},0) -- ++(-\tw,-\thR) -- ++(2*\tw,0) -- cycle;
			\draw ({2*\a-\gw},-\thR) -- ({2*\a+\gw},-\thR);
			\draw ({2*\a-\bw},-\thR-0.12) circle (2.4pt);
			\draw ({2*\a+\bw},-\thR-0.12) circle (2.4pt);
			\draw ({2*\a-\gw},-\thR-0.24) -- ({2*\a+\gw},-\thR-0.24);
			\foreach \x in {-0.24,-0.12,0,0.12,0.24}
			\draw[thin] ({2*\a+\x},-\thR-0.24) -- ++(-0.09,-0.11);
			\node[nodo] at ({2*\a},0) {};
			
			%% quotatura
			\draw[tacca] (0,-0.90) -- (0,-1.10);
			\draw[tacca] (\a,-0.90) -- (\a,-1.10);
			\draw[tacca] ({2*\a},-0.90) -- ({2*\a},-1.10);
			
			\draw[<->, thin, gray] (0,-1.00) -- (\a,-1.00);
			\node[above=0.04] at ({0.5*\a},-1.00) {$a$};
			\draw[<->, thin, gray] (\a,-1.00) -- ({2*\a},-1.00);
			\node[above=0.04] at ({1.5*\a},-1.00) {$a$};
			
		\end{tikzpicture}
		\caption{Trave su tre appoggi (\ESEref{ex:trave_tre_appoggi}):
			cerniera in $A=(0,0)$, carrello verticale intermedio in
			$B=(a,0)$, carrello verticale terminale in $C=(2a,0)$;
			polo $O\equiv A$. Il vincolo ridondante --- uno dei tre
			è in eccesso rispetto al minimo necessario per un unico
			corpo rigido --- rende il sistema iperstatico di grado 1.}
		\label{fig:trave_tre_appoggi}
	\end{figure}
